\chapter{总结与展望}\label{chap:conclusion}

\section{总结}

本文围绕混合式双轮足机器人的结构设计、动力学建模与轻量级控制方法开展研究，完成了从机械样机设计、控制模型建立到仿真和实物实验验证的完整流程。研究工作的主要结论如下。

在机械结构设计方面，本文根据双轮足机器人平衡移动、腿长调节、越障和嵌入式实时控制需求，完成了样机总体方案设计。样机采用躯干集中承载、闭环腿部机构支撑调节、足端轮组驱动移动的结构形式，并对腿部连杆、躯干结构、髋关节轴系、链传动和轮组系统进行了设计。通过关键尺寸分析、链传动校核以及腿部关节电机和轮毂电机选型，验证了样机结构方案的可实现性，为后续控制系统搭建提供了硬件基础。

在建模与力矩映射方面，本文将复杂闭环腿部机构等效为可伸缩虚拟腿，建立了机器人平面等效运动学与动力学模型，并推导了 LQR 控制器所需的线性状态空间模型。针对实际样机采用的混合式闭环腿部结构，进一步建立了虚拟腿任务空间到关节空间的力矩映射关系。

在控制方法与实验验证方面，本文提出了 LQR-VMC 分层轻量级控制系统。LQR 用于机体平衡、速度跟踪和腿摆控制量分配，VMC 用于腿长方向支撑与高度调节，滚转稳定环节通过左右腿长差修正改善侧向姿态。MuJoCo 仿真验证了控制器在原地平衡、速度跟踪、扰动恢复和腿长调节等工况下的闭环响应。物理样机实验进一步表明，所设计的控制系统能够在真实硬件上实现起立平衡、速度调节、腿长跟踪和滚转稳定，验证了本文结构方案、建模方法和控制策略的有效性。

\section{展望}

本文仍有若干不足需要在后续研究中完善。首先，本文建立的主动力学模型以平面运动为主，对滚转方向和复杂三维接触的描述仍较简化。后续可引入三维动力学模型、接触力估计和转向状态前馈，以提高机器人在高速转向、横坡和侧向冲击工况下的稳定性。其次，实物样机仍受到链传动回差、结构弹性、执行器响应滞后和轮地接触变化的影响，后续可从机械刚度、传动精度、驱动带宽和传感器融合精度等方面继续优化。

在控制算法方面，本文采用的 LQR-VMC 方法结构清晰、计算量较小，适合当前样机的嵌入式实时实现。但随着机器人运动速度、地形复杂度和任务要求的提高，后续可进一步引入更高级的控制方法。例如，模型预测控制（MPC）能够在滚动优化过程中显式考虑动力学约束、执行器限幅和接触约束，已有轮式双足和四足机器人研究表明，MPC 可用于提升全身姿态跟踪和动态运动能力\cite{yu2023hotwheel,dicarlo2018mpc}。此外，强化学习（RL）在腿式机器人运动控制中也表现出较强的策略搜索和复杂地形适应能力，相关研究已经实现了从仿真训练到实机部署的动态运动技能学习\cite{hwangbo2019learning,miki2022perceptive,rudin2021minutes}。因此，后续可以在现有轻量级控制框架的基础上，将 MPC 用于高层轨迹和接触力规划，或将 RL 用于参数自适应、复杂地形策略生成和鲁棒控制补偿，从而提升机器人在非结构化环境中的自主运动能力。
